\documentclass[12pt]{article}
\usepackage[utf8]{inputenc}
\usepackage[T1]{fontenc}
\usepackage[margin=1in]{geometry}
\usepackage{setspace}
\usepackage{xcolor}
\usepackage[colorlinks=true, urlcolor=blue!60!black, linkcolor=blue!60!black]{hyperref}
\setlength{\parindent}{0em}
\setlength{\parskip}{0.6em}
\renewcommand{\baselinestretch}{1.0}
\begin{document}
\begin{center}
{\LARGE \textit{Cooper v. Aaron}, 358 U.S. 1 (1958)\par}
\vspace{0.5em}
{\large Case Brief --- Two-Page Edition\par}
\end{center}

\section{Memory Jogger}
Official defiance cannot justify suspending desegregation, and the Court's interpretation of the Fourteenth Amendment binds every state official (\href{https://supreme.justia.com/cases/federal/us/358/1/}{Justia}).

\section{Detailed Case Facts}
The Little Rock School Board adopted a gradual desegregation plan beginning at Central High School in 1957. Arkansas pursued ``Massive Resistance,'' and Governor Orval Faubus used the National Guard to keep the ``Little Rock Nine'' out. After a mob forced the students' removal, President Eisenhower sent Army troops, and eight of the nine finished the year.

\section{Procedural History}
Citing ``chaos, bedlam and turmoil,'' the Board won a district court order suspending the plan for two and one-half years. The Eighth Circuit reversed. In a special term, the Supreme Court affirmed on September 12, 1958, and filed its opinion on September 29.

\section{Judicial Votes}
9--0, in an opinion personally signed by all nine Justices; Frankfurter also concurred.

\section{Holding}
School officials, as agents of the State, cannot suspend constitutional rights because of disorder provoked by the Governor and Legislature. Under \textit{Marbury} and Article VI, \textit{Brown}'s interpretation of the Fourteenth Amendment is the supreme law of the land, binding all state officials.

\section{Analysis of Opinions}
\textbf{Joint opinion:} The turmoil was ``directly traceable'' to state officials, and ``law and order are not here to be preserved by depriving the Negro children of their constitutional rights.'' Because the federal judiciary is supreme in expounding the Constitution, officials who swear to support it may not ``war against the Constitution.'' \textbf{Frankfurter:} yielding would ``enthrone official lawlessness''; orderly legal change differs from defiance. Later, \textit{United States v. Nixon} and \textit{City of Boerne} extended the principle to the federal branches.

\section{Examples: Future Applications}
\textbf{Same-Side:} a county clerk who stops issuing marriage licenses after the Court's ruling is bound and may be held in contempt. \textbf{Opposite-Side:} a President may veto a bill the Court would uphold. \textbf{Fence-Sitter:} a statute enacted to invite the Court to overrule a precedent.

\section{Critique}
Article VI makes the Constitution, not the Court's opinions, supreme, so critics such as Edwin Meese III say \textit{Cooper} confuses the two. Taken literally, it would condemn Lincoln's response to \textit{Dred Scott}.

\section{Key Quotations}
\begin{quote}
``This decision declared the basic principle that the federal judiciary is supreme in the exposition of the law of the Constitution, and that principle has ever since been respected by this Court and the Country as a permanent and indispensable feature of our constitutional system. It follows that the interpretation of the Fourteenth Amendment enunciated by this Court in the Brown case is the supreme law of the land, and Art.\ VI of the Constitution makes it of binding effect on the States `any Thing in the Constitution or Laws of any State to the Contrary notwithstanding.'\thinspace''
\end{quote}
\noindent --- Opinion of the Court, 358 U.S. at 18.

\end{document}
